% ============================================================
%  sections/02_related.tex
% ============================================================
\section{Background and Related Work}
\label{sec:related}

\subsection{Magnitude Pruning}

The idea of removing unimportant weights from neural networks dates to
\citet{lecun1990optimal}, who used second-order information (Hessian diagonal)
to identify low-salience weights.
\citet{han2015learning} introduced the modern magnitude-pruning pipeline:
train $\to$ prune low-$|w|$ weights $\to$ fine-tune.
\citet{han2015deep} extended this to a compression pipeline achieving
$35\times$ reduction on AlexNet.

\paragraph{Structured vs.\ unstructured pruning.}
\emph{Unstructured} pruning (our focus) zeros individual weights without
regard to their position, yielding sparse weight tensors.
While theoretically appealing, unstructured sparsity does not reduce
inference time on standard dense hardware without sparse-BLAS kernels or
specialized hardware~\citep{gale2019state}.
\emph{Structured} pruning removes entire channels or filters, making the
network genuinely smaller and faster~\citep{li2016pruning,molchanov2017pruning};
this is left as future work.

\citet{zhu2018prune} show that large sparse models consistently outperform
small dense models at equivalent parameter counts—suggesting that the
network topology matters as much as the weight values.

\subsection{Post-Training Quantization}

Quantization reduces the bit-width of weights and activations.
\citet{jacob2018quantization} formalize the INT8 affine quantization scheme
used in TensorFlow Lite and PyTorch's FBGEMM backend.
\citet{hubara2017quantized} study training with quantized representations.

\emph{Dynamic PTQ} (our primary method) quantizes only weights;
activations are quantized at runtime per-batch.
It requires no calibration data and is suitable for our setting.
\emph{Static PTQ} also quantizes activations using calibration statistics,
generally achieving better accuracy at the cost of requiring representative data.

On x86 CPUs, INT8 operations are accelerated by AVX2/AVX-512 SIMD instructions
via the FBGEMM backend, making quantization one of the few compression methods
where theoretical savings (4$\times$ fewer bytes) translates to real
latency reduction~\citep{deng2020model}.

\subsection{Low-Rank Factorization}

\citet{denton2014exploiting} demonstrate that convolutional filter responses
lie in low-dimensional subspaces, enabling $2\times$ speedup via SVD
decomposition of filter tensors.
\citet{li2018constrained} extend this to constrained low-rank training.
\citet{khodak2021initialization} study SVD-based initialization for
factorized architectures.

The key distinction between \emph{reconstruction}
(our approach: write back $\hat{\W} = \mathbf{U}_r \diag(\mathbf{s}_r) \mathbf{V}_r^\top$)
and \emph{factorization} (replace one Linear layer with two: $\R^n \to \R^r \to \R^m$)
is important: only the latter achieves actual parameter and memory reduction.

\subsection{Comparative Benchmarks}

\citet{blalock2020state} conduct a systematic review of pruning literature,
finding that comparisons across papers are often not reproducible due to
inconsistent evaluation protocols.
They call for standardized benchmarks—our work answers this call at
small scale.
\citet{gale2019state} show that simple magnitude pruning with careful
hyperparameter tuning is competitive with complex pruning algorithms.
Our work confirms this finding while adding quantization and low-rank as
comparators.
